\documentclass[10pt,a4paper]{amsart}
\usepackage{amsmath,amssymb,amsthm,graphicx}
\usepackage{verbatim}
\newenvironment{reference}{\comment}{\endcomment}
%\newenvironment{reference}{}{}
\newenvironment{slogan}{\comment}{\endcomment}
\newenvironment{history}{\comment}{\endcomment}

% For commutative diagrams we use Xy-pic
\usepackage[all]{xy}

% We use 2cell for 2-commutative diagrams.
\xyoption{2cell}
% Theorem environments.
%
\theoremstyle{plain}
\newtheorem{theorem}[subsection]{Theorem}
\newtheorem{proposition}[subsection]{Proposition}
\newtheorem{lemma}[subsection]{Lemma}

\theoremstyle{definition}
\newtheorem{definition}[subsection]{Definition}
\newtheorem{example}[subsection]{Example}
\newtheorem{exercise}[subsection]{Exercise}
\newtheorem{situation}[subsection]{Situation}

\theoremstyle{remark}
\newtheorem{remark}[subsection]{Remark}
\newtheorem{remarks}[subsection]{Remarks}

\numberwithin{equation}{subsection}

% Macros
%
\def\lim{\mathop{\mathrm{lim}}\nolimits}
\def\colim{\mathop{\mathrm{colim}}\nolimits}
\def\Spec{\mathop{\mathrm{Spec}}}
\def\Hom{\mathop{\mathrm{Hom}}\nolimits}
\def\Ext{\mathop{\mathrm{Ext}}\nolimits}
\def\SheafHom{\mathop{\mathcal{H}\!\mathit{om}}\nolimits}
\def\SheafExt{\mathop{\mathcal{E}\!\mathit{xt}}\nolimits}
\def\Sch{\mathit{Sch}}
\def\Mor{\mathop{\mathrm{Mor}}\nolimits}
\def\Ob{\mathop{\mathrm{Ob}}\nolimits}
\def\Sh{\mathop{\mathit{Sh}}\nolimits}
\def\NL{\mathop{N\!L}\nolimits}
\def\CH{\mathop{\mathrm{CH}}\nolimits}
\def\proetale{{pro\text{-}\acute{e}tale}}
\def\etale{{\acute{e}tale}}
\def\QCoh{\mathit{QCoh}}
\def\Ker{\mathop{\mathrm{Ker}}}
\def\Im{\mathop{\mathrm{Im}}}
\def\Coker{\mathop{\mathrm{Coker}}}
\def\Coim{\mathop{\mathrm{Coim}}}

% Boxtimes
%
\DeclareMathSymbol{\boxtimes}{\mathbin}{AMSa}{"02}

%
% Macros for moduli stacks/spaces
%
\def\QCohstack{\mathcal{QC}\!\mathit{oh}}
\def\Cohstack{\mathcal{C}\!\mathit{oh}}
\def\Spacesstack{\mathcal{S}\!\mathit{paces}}
\def\Quotfunctor{\mathrm{Quot}}
\def\Hilbfunctor{\mathrm{Hilb}}
\def\Curvesstack{\mathcal{C}\!\mathit{urves}}
\def\Polarizedstack{\mathcal{P}\!\mathit{olarized}}
\def\Complexesstack{\mathcal{C}\!\mathit{omplexes}}
% \Pic is the operator that assigns to X its picard group, usage \Pic(X)
% \Picardstack_{X/B} denotes the Picard stack of X over B
% \Picardfunctor_{X/B} denotes the Picard functor of X over B
\def\Pic{\mathop{\mathrm{Pic}}\nolimits}
\def\Picardstack{\mathcal{P}\!\mathit{ic}}
\def\Picardfunctor{\mathrm{Pic}}
\def\Deformationcategory{\mathcal{D}\!\mathit{ef}}
\usepackage{cleveref}
\crefname{originalreference}{original reference}{original references}
\Crefname{originalreference}{Original reference}{Original references}
\makeatletter
\newlabel{lemma-characterize-normal}{{09HQ}{source}{}{sourcepointer-1}{}}
\newlabel{lemma-characterize-normal@cref}{{[originalreference][0][]09HQ}{[1][0][]source}{}{}{}}
\makeatother
\crefname{originalreference}{original reference}{original references}
\Crefname{originalreference}{Original reference}{Original references}
\makeatletter
\newlabel{lemma-separable-degree}{{09HJ}{source}{}{sourcepointer-2}{}}
\newlabel{lemma-separable-degree@cref}{{[originalreference][0][]09HJ}{[1][0][]source}{}{}{}}
\makeatother
\title{The normal closure is a quotient of a tensor power of length the separable degree}
\author{The Stacks Project Authors}
\date{stacks-project:fields.tex:local-snapshot-20261012; GNU FDL 1.2 or later}
\begin{document}
\maketitle
\noindent\textit{Editorial note.} All delivered mathematical text is the author's own, copied from the pinned Stacks Project chapter file fields.tex; no statement or proof step was rewritten, reordered or removed.\par\medskip
\noindent\textit{Editorial note.} The counted proof is the author's complete argument: The tensor map is well defined because it is multilinear over K; surjectivity restates that M' is generated as a K-algebra by the sigma\_i(L).\par\medskip
\noindent\textit{Editorial note.} The label lemma-normal-closure-tensor-product names the counted statement, and the proof environment immediately following it in the source is delivered verbatim as its proof.\par\medskip
\noindent\textit{Editorial note.} The counted statement is carried by the author's own lemma environment in the chapter source.\par\medskip
\noindent\textit{Allowed standard prerequisites (editorial).}\begin{itemize}
\item Same chapter: lemma-separable-degree ([L : K]\_s = |Mor\_K(L, K-bar)|), lemma-characterize-normal (normality via equal images of embeddings, and its converse), uniqueness of the normal closure (lemma-normal-closure); standard fact: a K-subalgebra of an algebraic closure is a field (domains algebraic over a field are fields)
\end{itemize}\medskip
\ifdefined\hypertarget\hypertarget{sourcepointer-1}{}\fi
\noindent\textit{Source reference.} Original reference 09HQ: sources/stacks-fields/pinned-original-source.tex, source lines 2096--2096. Prerequisite pointer invoked at 2403; it names a standard fact of the same chapter that the counted argument uses, and every step of the counted proof built on it is retained verbatim in the delivered body.\par\medskip
\ifdefined\hypertarget\hypertarget{sourcepointer-2}{}\fi
\noindent\textit{Source reference.} Original reference 09HJ: sources/stacks-fields/pinned-original-source.tex, source lines 1968--1968. Prerequisite pointer invoked at 2395; it names a standard fact of the same chapter that the counted argument uses, and every step of the counted proof built on it is retained verbatim in the delivered body.\par\medskip
\begin{lemma}
\label{lemma-normal-closure-tensor-product}
Let $L/K$ be a finite extension. Let $M/L$ be the normal
closure of $L$ over $K$. Then there is a surjective map
$$
L \otimes_K L \otimes_K \ldots \otimes_K L \longrightarrow M
$$
of $K$-algebras where the number of tensors can be taken
$[L : K]_s \leq [L : K]$.
\end{lemma}

\begin{proof}
Choose an algebraic closure $\overline{K}$ of $K$.
Set $n = [L : K]_s = |\Mor_K(L, \overline{K})|$ with equality by
Lemma \ref{lemma-separable-degree}.
Say $\Mor_K(L, \overline{K}) = \{\sigma_1, \ldots, \sigma_n\}$.
Let $M' \subset \overline{K}$ be the $K$-subalgebra
generated by $\sigma_i(L)$, $i = 1, \ldots, n$. Then $M'$ is
a field since any $K$-subalgebra of $\overline{K}$ is a field.
Any $K$-algebra map from $M'$ to $\overline{K}$ permutes the $\sigma_i$
so sends $M'$ into and onto $M'$. By construction the field $M'$
is generated by conjugates of elements of $\sigma_1(L)$. Having said
this it follows from Lemma \ref{lemma-characterize-normal}
that $M'$ is normal over $K$ and that it is the
smallest normal subextension of $\overline{K}$ containing
$\sigma_1(L)$. By uniqueness of normal closure we have $M \cong M'$.
Finally, there is a surjective map
$$
L \otimes_K L \otimes_K \ldots \otimes_K L \longrightarrow M',
\quad
\lambda_1 \otimes \ldots \otimes \lambda_n \longmapsto
\sigma_1(\lambda_1) \ldots \sigma_n(\lambda_n)
$$
and note that $n \leq [L : K]$ by definition.
\end{proof}

\end{document}
